\documentclass[12pt]{article}
\usepackage[utf8]{inputenc}
\usepackage[spanish]{babel}
\usepackage{hyperref}
\usepackage{graphicx}
\usepackage{listings}
\usepackage{geometry}
\geometry{a4paper, margin=1in}

\title{Antigravity CLI (AGY) for Dummies\\
\large La guía amigable para usar a tu nuevo asistente virtual}
\author{Aldo Zetina}
\date{\today}

\begin{document}

\maketitle

\begin{abstract}
¿Solo usas Excel, Word y Redes Sociales? ¡Perfecto! Este manual está hecho a tu medida. Aquí aprenderás a usar a tu nuevo asistente virtual (una Inteligencia Artificial) sin usar palabras complicadas. Imagina que es como chatear por WhatsApp con un becario súper rápido y muy inteligente que está listo para hacer tu trabajo aburrido.
\end{abstract}

\tableofcontents
\newpage

\section{Conociendo a tu Asistente}
\subsection{¿Qué es una Inteligencia Artificial (IA)?}
No te compliques: piénsalo como un asistente personal muy aplicado. No es un robot de película, es solo un programa con el que puedes "chatear" para pedirle que te ayude con tareas de tu día a día, como organizar información, resumir documentos o redactar correos.

\subsection{¿Qué es esa "pantalla negra" (la terminal)?}
Para hablar con este asistente usamos algo llamado "terminal". Olvídate de los botones y el ratón; la terminal es simplemente una ventana de chat donde tú escribes un mensaje, presionas Enter, y el asistente te responde. ¡Así de fácil!

\subsection{Cómo abrir la terminal en Windows y saludar}
\begin{enumerate}
    \item Presiona la tecla \textbf{Windows} en tu teclado (la que tiene el logo de la ventanita).
    \item Escribe la palabra \texttt{cmd} y presiona \textbf{Enter}.
    \item Se abrirá una ventana negra. Ahí, simplemente escribe la palabra \texttt{agy} (que es el nombre de tu asistente) y presiona \textbf{Enter}.
\end{enumerate}
¡Listo! La pantalla cambiará y verás que tu asistente te saluda. Ya están listos para trabajar. Cuando termines, solo escribe \texttt{/exit} y presiona Enter para despedirte.

\section{El Arte de Pedir (Prompts)}
\subsection{¿Qué es un "prompt"?}
Es solo una palabra de moda para referirse a la \textbf{instrucción o petición} que le escribes a tu asistente. Es literalmente el mensaje que le envías en el chat.

\subsection{La receta para pedir cosas bien}
Tu asistente es muy inteligente, pero no lee la mente. Si le das instrucciones vagas, el resultado no te gustará. Un buen mensaje (prompt) siempre tiene tres partes:
\begin{enumerate}
    \item \textbf{El Rol:} Dile cómo quieres que actúe (ej. "Actúa como un redactor experto" o "Como un contador").
    \item \textbf{La Tarea:} Qué quieres que haga (ej. "Redacta un correo", "Resume este archivo").
    \item \textbf{Los Detalles:} Para quién es, qué tono usar (formal o amigable), qué no debe incluir.
\end{enumerate}

\subsection{Ejemplos: De malo a excelente}
Veamos cómo mejora una instrucción para redactar un correo:
\begin{itemize}
    \item \textbf{Mensaje malo:} "Escribe un correo pidiendo vacaciones." (El asistente no sabe a quién va dirigido ni las fechas).
    \item \textbf{Mensaje regular:} "Escribe un correo para mi jefe Carlos pidiendo vacaciones para la próxima semana." (Mejor, pero le falta tono).
    \item \textbf{Mensaje Excelente:} "Actúa como un empleado muy respetuoso. Escribe un correo corto y formal para mi jefe Carlos, solicitando mis vacaciones del 10 al 15 de agosto. Menciona que dejaré todos mis pendientes terminados."
\end{itemize}

\section{Ejemplo Práctico: Extracción y Manejo de Datos (PDF a Excel)}
\subsection{Tu asistente como administrador de datos}
Aquí viene lo mejor. Tu asistente no solo redacta correos; ¡puede hacer el trabajo pesado de administración por ti! Es capaz de leer documentos grandísimos (como un PDF de 100 páginas) y extraer solo la información que te interesa, poniéndola directamente en un archivo de Excel, sin que tú tengas que copiar y pegar celda por celda.

\subsection{Qué tipo de "magia" puede hacer}
\begin{itemize}
    \item \textbf{De PDF a Excel:} "Saca todas las tablas de este reporte en PDF y pásalas a un Excel."
    \item \textbf{Limpieza de datos:} "Lee mi lista de Excel y borra todas las filas repetidas."
    \item \textbf{Clasificación automática:} "Toma esta lista de 100 productos y agrúpales una categoría."
\end{itemize}

\subsection{Cómo pedirle que trabaje con Excel}
\textbf{Ejemplo de Prompt Excelente para Excel:}\\
\textit{"Actúa como mi auxiliar administrativo. Lee el archivo 'facturas.pdf' que tengo guardado. Busca todas las compras del mes de marzo. Saca solamente las columnas de 'Fecha', 'Producto' y 'Total', y guarda todo ordenado en un nuevo archivo de Excel llamado 'compras\_marzo.xlsx'."}

\section{Los "Atajos" de tu Asistente}
El asistente tiene unos comandos rápidos (siempre empiezan con una barra diagonal \texttt{/}) para facilitarte la vida:
\begin{itemize}
    \item \textbf{\texttt{/plan}}: Úsalo cuando la tarea es muy grande. Le dice al asistente: "Antes de hacer nada, dime paso a paso cómo lo vas a hacer para que yo te lo apruebe".
    \item \textbf{\texttt{/goal}}: Le dice al asistente que se concentre totalmente en la tarea hasta terminarla, ideal para cuando le pides algo que va a tardar mucho.
\end{itemize}

\section{¿Y si se equivoca?}
No te preocupes ni te pelees con la computadora. Si el asistente te entrega algo que no te gusta, simplemente respóndele en el chat explicándole qué hizo mal. Por ejemplo: \textit{"Ese correo suena muy formal, hazlo más amigable"} o \textit{"Te faltó incluir la columna de precios en el Excel, agrégala por favor"}. Trátalo como a una persona y te entenderá perfecto.

\end{document}
